% ── Related Work ──────────────────────────────────────────────────────────────
\section{Related Work}

\subsection{Dense Retrieval and RAG}

Classical retrieval relies on lexical overlap, epitomised by Okapi BM25
\parencite{robertson2009bm25}, which remains a strong and ubiquitous baseline.
Dense retrieval instead embeds queries and passages into a shared vector space
so that semantically related text is close even without shared terms;
\textcite{karpukhin2020dpr} demonstrated that such dense passage retrieval can
substantially outperform BM25 on open-domain question answering. Sentence-level
embedding models \parencite{reimers2019sbert} made this approach practical at
scale. \textcite{lewis2020rag} introduced retrieval-augmented generation, coupling
a dense retriever with a generative model so that answers are grounded in
retrieved evidence rather than parametric memory alone. This project builds a
Granite-based instantiation of this retrieve-then-generate paradigm.

\subsection{Evaluation Benchmarks and Metrics}

Rigorous retrieval evaluation requires datasets with human relevance judgments.
MS~MARCO \parencite{nguyen2016msmarco} provides large-scale passage-ranking data
from real search queries; Natural Questions \parencite{kwiatkowski2019nq} offers
open-domain questions answered from Wikipedia; and BEIR
\parencite{thakur2021beir} aggregates heterogeneous retrieval tasks into a single
zero-shot benchmark suite, making it the de-facto standard for comparing
retrieval systems. Performance is reported with standard
information-retrieval metrics---precision@k, recall@k, nDCG, and MRR---which this
project adopts for the system-versus-baseline comparison.

\subsection{Long-Context Retrieval (Secondary)}

As a complementary diagnostic, \textcite{liu2024lost} show that language models
exhibit a U-shaped retrieval curve over long inputs---accuracy is highest when
relevant text sits at the start or end and collapses in the middle (the
\emph{Lost-in-the-Middle} effect). This motivates the project's secondary
long-context stress test.

% ── System Design ─────────────────────────────────────────────────────────────
\section{System Design}

\subsection{Architecture Overview}

The system has four stages: an \emph{ingestion} pipeline that loads a corpus,
chunks it, and builds a persistent vector index of Granite embeddings; a
\emph{retrieval} layer offering both the Granite-embedding dense retriever (the
delivered system) and a BM25 baseline; a
\emph{RAG} layer that generates answers grounded in the retrieved passages; and
an \emph{explainability} layer that attributes each answer to its supporting
sources.

\subsection{Technology Stack}

\begin{itemize}
  \item \textbf{IBM watsonx.ai (Granite)} --- the \texttt{granite-embedding}
        models provide the query and passage embeddings that \emph{define} the
        dense retriever (the component whose retrieval quality this project
        measures), and a Granite generative model produces the grounded answers
        in the RAG layer. Using Granite embeddings for retrieval itself---rather
        than only for generation---is what makes the system a genuinely
        \emph{Granite-powered retrieval} system and gives the
        Granite-vs-baseline hypothesis its meaning.
  \item \textbf{LangChain / \texttt{langchain-ibm}} --- orchestration and
        watsonx.ai integration.
  \item \textbf{FAISS} --- persistent dense vector index.
  \item \textbf{sentence-transformers} --- open-source embedding baseline.
  \item \textbf{rank\_bm25} --- classical BM25 baseline.
  \item \textbf{BEIR / Hugging Face \texttt{datasets}} --- standard benchmark data.
  \item \textbf{ranx / pytrec\_eval} --- precision/recall, nDCG, MRR.
  \item \textbf{Streamlit} --- interactive demo with cited answers.
\end{itemize}

\subsection{Explainability and Trust}

To address the brief's requirement for trust in retrieved outputs, every
generated answer is attributed to the specific chunks that support it
(\texttt{src/explainability}), so users can trace each claim back to its source
passage and page.

% ── Evaluation Plan ───────────────────────────────────────────────────────────
\section{Evaluation Plan}

\subsection{Primary: Benchmark Comparison}

The headline experiment indexes a benchmark corpus once and runs three
retrievers over the same queries: the Granite dense retriever (the delivered
system), the BM25 baseline \parencite{robertson2009bm25}, and an open-source
dense retriever \parencite{karpukhin2020dpr}. Each run is scored against the
benchmark's relevance judgments using:

\begin{description}
  \item[Precision@k] of the top-$k$ retrieved passages, the fraction relevant.
  \item[Recall@k] of all relevant passages, the fraction appearing in the top-$k$.
  \item[nDCG@k] rank-aware quality, rewarding relevant hits near the top.
  \item[MRR] mean reciprocal rank of the first relevant hit.
\end{description}

Results are reported as a system-versus-baseline comparison table across
$k \in \{1,3,5,10\}$, testing the hypothesis that the Granite system improves
precision and recall over both baselines.

\subsection{Secondary: RAG Answer Quality and Long-Context Stress Test}

The RAG layer's answers are additionally assessed with RAGAS
\parencite{es2023ragas} for context precision and faithfulness. Separately, a
long-context ``needle'' stress test injects a known fact at controlled depths in
long documents and measures retrieval accuracy across document length and
injection depth, quantifying the \emph{Lost-in-the-Middle} effect
\parencite{liu2024lost} for the system.

% ── Time Plan ─────────────────────────────────────────────────────────────────
\section{Time Plan}

\begin{center}
\renewcommand{\arraystretch}{1.35}
\begin{tabular}{@{} >{\bfseries}p{3.2cm} p{5.8cm} p{5cm} @{}}
\toprule
\textbf{Period} & \textbf{Milestone} & \textbf{Deliverable} \\
\midrule
Wk 1--2 \newline (Jun 9--20) &
  Environment setup; watsonx.ai + baseline client; load a standard benchmark
  (corpus/queries/qrels). &
  API smoke test passes; one benchmark loads end-to-end. \\

Wk 3--4 \newline (Jun 23 -- Jul 4) &
  Ingestion pipeline (chunk, embed, persistent FAISS index); BM25 baseline. &
  Indexed corpus; working BM25 retriever. \\

Wk 5--6 \newline (Jul 7--18) &
  Granite dense retriever; implement precision/recall, nDCG, MRR. &
  First system-vs-BM25 metrics on one benchmark. \\

Wk 7--8 \newline (Jul 21 -- Aug 1) &
  RAG generation layer; explainability / citations. &
  End-to-end cited answers; RAGAS scores. \\

Wk 9--10 \newline (Aug 4--15) &
  Full benchmark sweep across datasets; add open-source dense baseline. &
  System-vs-baselines comparison tables and figures. \\

Wk 11--12 \newline (Aug 18--29) &
  Streamlit demo (upload, search, cited answers); optional long-context
  needle stress test. &
  Live demo; secondary stress-test results. \\

Wk 13 \newline (Sep 1--4) &
  Write individual report; final code review and repository clean-up. &
  \textbf{Final submission by 4 September 2026.} \\
\bottomrule
\end{tabular}
\end{center}
